\begin{textsample}{POS Dim 1 – ma – Score 42.00 – ma/ma\_em\_24.txt}  \label{ex:f1_pos_066}
1.7. \textbf{Uma} introdução às áreas de conhecimentos De acordo com a proposta da BNCC, a organização das aprendizagens está dividida em quatro áreas de conhecimento, a saber : Ciências Humanas e Sociais Aplicadas ; Linguagens e suas Tecnologias ; Matemática e Ciências da Natureza.
A área de Ciências Humanas e Sociais Aplicadas para o ensino médio tem por finalidade a ampliação e o aprofundamento das aprendizagens essenciais desenvolvidas no ensino fundamental.
É importante observar \textbf{que}, no ensino fundamental, a área era composta pelos componentes curriculares geografia e história e \textbf{que}, a partir do ensino médio, a composição dessa área passa a reunir, além de história e geografia, os componentes filosofia e sociologia.
A área tem as aprendizagens nos ensinos fundamental e médio \textbf{baseadas} no desenvolvimento de competências \textbf{que} visam à identificação, análise, comparação e interpretação das ideias, pensamentos, fenômenos e processos históricos, geográficos, sociais, econômicos, políticos e culturais.
Assim sendo, no ensino médio, a área direciona as aprendizagens de forma a aprofundá -las com expansão da base \textbf{conceitual} e dos modos de construção da argumentação e \textbf{sistematização} do \textbf{raciocínio}, sob a perspectiva de procedimentos analíticos e interpretativos, fortalecidos pela inclusão da filosofia e da sociologia.
A área de Ciências Humanas e Sociais Aplicadas tem por objeto de estudo o \textbf{sujeito} e suas relações construídas pelos diferentes indivíduos, grupos, segmentos e classes sociais.
Analisar o \textbf{sujeito} e o \textbf{viver} em sociedade faz da área \textbf{uma} importante aliada no processo de construção de conhecimento para a convivência social, a partir da construção de \textbf{uma} articulação entre os componentes filosofia, geografia, história e sociologia, permitindo o desenvolvimento de competências e habilidades importantes para a compreensão e o reconhecimento das diferenças, o respeito aos direitos humanos e o combate à intolerância e preconceitos de qualquer natureza.
Outro aspecto importante para a área refere -se à \textbf{divisão} por \textbf{categorias} temáticas \textbf{que} a BNCC aponta como prioritárias para o desenvolvimento das aprendizagens essenciais e \textbf{que} englobam a filosofia, a geografia, a história e a sociologia, de forma a garantir \textbf{um} currículo \textbf{que} \textbf{fomenta} a ação interdisciplinar, no sentido de \textbf{que} ambos os componentes se associem na construção de conceitos cada vez mais complexos.
Nessa configuração, as \textbf{categorias} Tempo e Espaço ; Território e \textbf{Fronteira} ; Indivíduo, Natureza, Sociedade, Cultura e Ética ; e Política e Trabalho aprofundam as \textbf{abordagens} na busca do entendimento sobre o porquê da nossa constituição social e das problemáticas \textbf{contemporâneas} \textbf{que} acompanham a nossa sociedade \textbf{atual}.
Na \textbf{categoria} Tempo e Espaço, por exemplo, é fundamental a compreensão das diferentes noções de tempo para entender as diferenças de determinadas sociedades e culturas.
Ao analisarmos a \textbf{categoria} Território e \textbf{Fronteira}, a temática sobre globalização e limites das \textbf{fronteiras} das nações se faz necessária, pois explica muitos problemas da \textbf{contemporaneidade}.
Com relação à \textbf{categoria} Indivíduo, Natureza, Sociedade, Cultura e Ética, temos \textbf{aqui} a base para o desenvolvimento de importantes concepções para a área, como a busca pelo desenvolvimento de \textbf{um} jovem com valores voltados ao respeito, à convivência, ao bem comum e à ética.
Outra \textbf{categoria} importante é a de Política e Trabalho, \textbf{que} promove o desenvolvimento do jovem cidadão dotado de senso crítico e entende \textbf{que}, a esse/a jovem, deve ser dado o direito ao conhecimento do mundo político e econômico em \textbf{que} \textbf{vive}, levando em consideração os temas \textbf{contemporâneos} como elementos do processo ensino e aprendizagem de grande relevância.
Segundo Edgar Morin ( 2015 ), é \textbf{preciso} ensinar a compreensão humana, e as ciências humanas têm essa difícil \textbf{tarefa} de ajudar o \textbf{educando} a descortinar o mundo por meio de \textbf{uma} reflexão crítica \textbf{que} contemple a ética, os direitos humanos e o autoconhecimento, considerando competências, habilidades e suas relações com o trabalho das \textbf{humanidades}.
Na Base Nacional Comum Curricular ( BNCC ), a área de Linguagens e suas Tecnologias visa \textbf{aprimorar}, consolidar e ampliar as aprendizagens \textbf{vistas} durante o ensino fundamental desses componentes curriculares.
Conforme consta na BNCC ( 2017 ), a área de Linguagens e suas Tecnologias no ensino médio vem para ajudar o estudante a ampliar sua autonomia, seu protagonismo e sua autoria nas práticas de diferentes linguagens, proporcionando \textbf{uma} visão crítica aos diferentes usos das linguagens na argumentação, na apreciação e na participação das diferentes possibilidades de manifestações artísticas e culturais e, atrelado a isso, o fomento ao uso criativo das diversas mídias.
Além disso, essa área prioriza cinco campos de atuação social : o campo da vida pessoal ; o campo das práticas de estudo e pesquisa ; o campo jornalístico-midiático ; o campo de atuação na vida pública ; e o campo artístico.
No ensino médio, nos diferentes componentes curriculares da área de Linguagens e suas Tecnologias, procura -se garantir aos estudantes aprofundar aprendizagens, consolidar autonomia, promover valores e desenvolver habilidades, buscando \textbf{uma} formação integral \textbf{que} amplie as práticas de linguagem e dos repertórios, a diversificação dos campos nos quais atuem, a análise das manifestações artísticas, corporais e linguísticas e de como essas manifestações constituem a vida social em diferentes culturas pertencentes a variadas localidades do nosso território nacional.
Reiteramos \textbf{que} o desafio da área de Linguagens e suas Tecnologias no ensino médio é consolidar e garantir a continuidade do \textbf{que} foi aprendido na etapa anterior ( ensino fundamental ), na perspectiva de superação de dificuldades no processo de aprendizagem \textbf{que}, eventualmente, se tenham acumulado.
Para tanto, é necessário valorizar os conhecimentos adquiridos e consolidados pelos estudantes, a fim de melhor planejar as intervenções pedagógicas adequadas para assegurar a eles \textbf{que} continuem aprendendo, visando, ainda, \textbf{um} nível de complexidade maior.
Portanto, destacarmos \textbf{que} os componentes curriculares da área de Linguagens e suas Tecnologias devem : a ) proporcionar \textbf{que} os estudantes trabalhem com a semiose, ou seja, \textbf{que} explorem as possibilidades \textbf{expressivas} das diversas linguagens ; b ) possibilitar \textbf{que} os estudantes compreendam e reflitam sobre a emancipação ; c ) garantir a inclusão na aprendizagem ; d ) permitir meios para \textbf{fomentar} o protagonismo juvenil ; d ) relacionar os conhecimentos específicos de cada componente com os Temas Transversais ; e ) considerar como base, na prática de ensino, o tratamento \textbf{metodológico} da \textbf{interdisciplinaridade} e \textbf{transdisciplinaridade}.
Sobre a área de Matemática, a BNCC propõe \textbf{que}, nessa etapa da educação básica, os estudantes desenvolvam \textbf{uma} visão mais integrada do componente curricular matemática, considerando, assim, a realidade e sua aplicação em diferentes contextos, devendo -se explorar o caráter instrumental dessa área do conhecimento.
Para isso, deve -se proporcionar \textbf{uma} ampliação e o aprofundamento das aprendizagens desenvolvidas ao longo do ensino fundamental.
Isso se dá por meio do desenvolvimento de habilidades específicas relacionadas a raciocinar, representar, comunicar e \textbf{argumentar}.
A ampliação e o aprofundamento são caminhos permanentes a serem seguidos, assim como no ensino fundamental, dando significação aos conhecimentos e considerando as experiências dos estudantes, incentivando a pesquisa e o entendimento dos fenômenos tecnológicos, científicos e sociais.
Além disso, \textbf{que} os preparem para o mercado de trabalho, desenvolvendo suas capacidades críticas e reflexivas.
A matemática deve ser ferramenta de construção, investigação e aplicação do conhecimento, possibilitando ao estudante intervir na sua realidade social de maneira ética e responsável.
Para ampliação e aprofundamento do conhecimento matemático, o estudante deve ser motivado a questionar, formular, testar e, posteriormente, validar as suas próprias hipóteses, verificando a adequação de sua resposta à situação-problema proposta, a partir da construção de formas de pensar \textbf{que} o levem a refletir e agir criticamente sobre as mais diversas questões cotidianas.
Isto se dá a partir de \textbf{um} processo de \textbf{ensino-aprendizagem} da matemática \textbf{que} vise a \textbf{uma} compreensão abrangente de mundo e \textbf{que} também \textbf{qualifique} a inserção do estudante no mundo do trabalho, capacitando -o para tornar sua argumentação consistente e dando -lhe segurança para lidar com problemas e desafios de origens diversas, de forma contextualizada e interdisciplinar, permitindo -lhe, ao mesmo, fazer uso de sua imaginação e criatividade nos diversos contextos.
Dessa forma, assim como no ensino fundamental, os temas \textbf{contemporâneos} transversais – \textbf{que} são temas \textbf{contemporâneos} \textbf{que} afetam a vida humana em escala global, regional e local – devem ser acolhidos pelo componente curricular de matemática no intuito de \textbf{fomentar} o \textbf{compromisso} do professor com a formação dos estudantes, pois os objetos de conhecimento devem ser trabalhados de forma a \textbf{explicitarem} as questões transversais \textbf{que} perpassam as várias \textbf{instâncias} do componente curricular, bem como toda a escola, em seu contexto social.
Assim, a matemática, tanto como área de conhecimento quanto como componente curricular, relaciona -se com os temas integradores, provendo \textbf{um} maior contato do estudante com o meio externo, interferindo na mudança dos valores e desenvolvendo o senso crítico e o \textbf{posicionamento} acerca das questões sociais, contribuindo com sua formação como cidadão.
Assim, os temas integradores surgem como complementos importantes aos procedimentos educacionais, \textbf{aproximando} o estudante não apenas do saber escolar, mas do saber matemático inserido nas mais diversas questões da sociedade.
O componente curricular matemática articula os conhecimentos específicos da área com as demais áreas do conhecimento por meio de seus objetos de conhecimentos, \textbf{que} serão apresentamos nos parágrafos seguintes.
As questões globais e locais com as quais a área de Ciência e Tecnologia está envolvida, e sua \textbf{influência} no modo como \textbf{vivemos}, pensamos, agimos e as consequências das ações antropogênicas passaram a incorporar as preocupações de muitos brasileiros.
Nesse contexto, a área de Ciências da Natureza tende a ser encarada não somente como ferramenta para solucionar problemas, mas também como \textbf{uma} abertura para novas visões de mundo.
Nossos aprendizados envolvem tanto a comunidade \textbf{que} nos é mais próxima quanto a cósmica.
Temos como meta ajudar a construir valores sociais voltados não apenas para a conservação do meio ambiente e sua sustentabilidade, mas também para valores críticos \textbf{que} se responsabilizem pelas modificações \textbf{que} ocorrem no ambiente natural, com a expectativa de \textbf{que} contribuam para \textbf{uma} melhor e mais sadia qualidade de vida.
Todavia, poucas pessoas aplicam os conhecimentos e procedimentos científicos na resolução de seus problemas cotidianos.
Tal constatação corrobora a necessidade de a área de Ciências da Natureza comprometer -se com o letramento científico da população.
Além da investigação científica, \textbf{uma} \textbf{abordagem} fundamental para a formação da cidadania é a \textbf{contextualização} \textbf{sócio-histórica} dos conteúdos das Ciências da Natureza, \textbf{que} não se resume à exemplificação de fatos ou situações cotidianas, mas “ [... ] abarca competências de inserção da ciência e de suas tecnologias em \textbf{um} processo histórico, social e cultural e o reconhecimento e discussão de aspectos práticos e éticos da ciência no mundo \textbf{contemporâneo} ” ( BRASIL, 2002, p. 31 ).
Nessa perspectiva, a BNCC da área, por meio de \textbf{um} olhar articulado da biologia, física e química, define as competências e habilidades no \textbf{que} se refere : aos conhecimentos \textbf{conceituais} da área ; à \textbf{contextualização} social, cultural, ambiental e histórica desses conhecimentos ; aos processos e práticas de investigação ; e às linguagens das Ciências da Natureza.
Na definição das competências e habilidades específicas da área, a BNCC propõe \textbf{um} aprofundamento nas temáticas para enfrentar desafios locais e/ou globais relativos às condições de vida e ao ambiente, dividido em três eixos temáticos : matéria e energia, vida e evolução e tecnologia e linguagem científica.
Essas temáticas constituem \textbf{uma} base \textbf{que} permite investigar, analisar e discutir situações-problema, além de compreender e interpretar leis, \textbf{teorias} e modelos, aplicando -os na resolução de problemas.
Dessa forma, os estudantes podem reelaborar seus próprios saberes e reconhecer as potencialidades e limitações das ciências ; aprofundar e ampliar suas reflexões a respeito dos contextos de produção e aplicação do conhecimento científico e tecnológico ; avaliar o impacto de tecnologias \textbf{contemporâneas} em seu cotidiano, em setores produtivos, na economia e nas dinâmicas sociais ; compreender a dinâmica da construção do conhecimento científico ; promover o protagonismo na aplicação dos \textbf{métodos} e procedimentos científicos ; contextualizar desafios e situações-problema como gatilho mental da curiosidade e criatividade, na elaboração de procedimentos e busca de soluções.
Isso permite \textbf{que} o diálogo com o mundo real se intensifique, bem como as possibilidades de análises e de intervenções em contextos mais amplos e complexos.
Vale destacar \textbf{que} o pensamento científico introduzido o mais cedo possível, a partir da educação infantil, por meio de atividades práticas como a observação do contexto e reflexões sobre os fatos sociais e fenômenos ambientais, irá possibilitar a ampliação da compreensão de seu contexto, de sua realidade e do mundo, o \textbf{que} contribuirá para o desenvolvimento dos conhecimentos das diferentes áreas.

% matched lemmas: abordagem, aprimorar, aproximar, aqui, argumentar, atual, basear, categoria, compromisso, conceitual, contemporaneidade, contemporâneo, contextualização, divisão, educar, ensino-aprendizagem, explicitar, expressivo, fomentar, fronteira, humanidade, influência, instância, interdisciplinaridade, metodológico, método, posicionamento, preciso, qualificar, que, raciocínio, sistematização, sujeito, sócio-histórico, tarefa, teoria, transdisciplinaridade, um, ver, vivar, viver
\end{textsample}
